%% kapa: exam=1 round=36 year=2025 period=1 form=A subject=economics subject_ko=경제학원론 seq=18 number=58
%% source: Q-Net 감정평가사 기출문제 (cst00309) artlSeq=5231525 file=raw/1008/5231525/2025년 제36회 감정평가사 1차시험 1교시 문제지 원본.pdf
%% conversion: kordoc 4.24.51 -> tools/gen1.py (auto); review=proofread-sonnet; flags=none
\begin{question}{58}
어떤 재화의 시장수요함수는 $Q=100-P$ 이고, 시장공급함수는 $Q=-10+P$ 이다. 정부가 $P=40$에서 가격상한제를 실시할 경우, 생산자 잉여의 감소분 중 소비자 잉여로 이전(transfer)되는 크기는? (단, $P$는 가격, $Q$는 수량)
\begin{choices}
  \item 350
  \item 400
  \item 450
  \item 500
  \item 550
\end{choices}
\end{question}
\answer{3}
